\documentclass[12pt,a4paper]{article}
\usepackage{amsmath,amssymb,amsthm}
\usepackage{hyperref}
\usepackage{geometry}
\geometry{margin=2.5cm}
\usepackage{enumitem}
\usepackage{booktabs}

\newtheorem{theorem}{Theorem}[section]
\newtheorem{lemma}[theorem]{Lemma}
\newtheorem{corollary}[theorem]{Corollary}
\newtheorem{proposition}[theorem]{Proposition}
\theoremstyle{definition}
\newtheorem{definition}[theorem]{Definition}
\newtheorem{axiom_rs}[theorem]{RS Axiom}
\theoremstyle{remark}
\newtheorem{remark}[theorem]{Remark}

\title{Special Relativity and Local Lorentz Invariance\\
from the Recognition Science Causal Bound}

\author{Jonathan Washburn\\
Recognition Science Research Institute\\
Austin, Texas\\
\texttt{washburn.jonathan@gmail.com}}

\date{March 2026}

\begin{document}
\maketitle

\begin{abstract}
Special Relativity (SR) --- the invariance of physical laws under Lorentz
transformations and the universal speed limit $c$ --- is conventionally
postulated.  We derive SR from the Recognition Science (RS) framework, in
which a single primitive, the Recognition Composition Law (RCL)
$J(xy)+J(x/y) = 2J(x)J(y)+2J(x)+2J(y)$, forces a discrete eight-tick
clock with fundamental tick $\tau_0$ and voxel size $\ell_0$.  The causal
speed bound $c = \ell_0/\tau_0$ emerges as the maximum rate at which
ledger information can propagate in one tick (Lean: \texttt{StepBounds},
proved).  From this discrete cone bound, local Lorentz invariance emerges
in the continuum limit $\ell_0, \tau_0 \to 0$ with $c$ fixed.  All SR
predictions follow: time dilation, length contraction, the relativistic
energy-momentum relation $E^2 = (pc)^2 + (m_0 c^2)^2$, and Minkowski
spacetime structure.  The RS derivation requires no postulate about the
isotropy of light; isotropy follows from the voxel lattice symmetry
$\mathbb{Z}^3$ (proved in \texttt{Verification.Necessity.VoxelSymmetry}).
The Lean certificate \texttt{ConeBoundCert} (proved, zero sorry) certifies
the discrete causal bound.  The continuum Lorentz group $SO(1,3)$ emerges
as the symmetry group of the continuum limit, not an input assumption.
\end{abstract}

%%============================================================
\section{Introduction}
\label{sec:intro}
%%============================================================

Einstein's two postulates of Special Relativity~\cite{Einstein1905} are:
\begin{enumerate}
  \item The laws of physics are the same in all inertial frames.
  \item The speed of light in vacuum is $c \approx 2.998 \times 10^8$
        m/s, the same for all observers regardless of source motion.
\end{enumerate}

From these postulates follow the Lorentz transformations, time dilation,
length contraction, and the equivalence of mass and energy.  Despite
immense experimental confirmation~\cite{Zhang1997}, the postulates
themselves have no deeper justification in standard physics.

In Recognition Science (RS)~\cite{Washburn2026a}, neither postulate is
assumed.  Both emerge from a single, deeper structure: the discrete
ledger dynamics on a $D=3$ voxel lattice with eight-tick period.  The
speed $c$ is derived as $c = \ell_0/\tau_0$, the ratio of the voxel
spacing to the atomic tick.  Lorentz invariance is the continuum-limit
symmetry of this discrete structure.

This paper is organized as follows.  Section~\ref{sec:rs} reviews the
RS elements used.  Section~\ref{sec:derivation} derives $c$ and the
causal cone.  Section~\ref{sec:SR} derives the full SR kinematic
structure.  Section~\ref{sec:lean} catalogs the Lean formalization.
Section~\ref{sec:discussion} discusses the RS predictions.

%%============================================================
\section{Recognition Science Framework}
\label{sec:rs}
%%============================================================

\subsection{The Recognition Composition Law}

\begin{axiom_rs}[RCL --- the unique RS primitive]
\begin{equation}
  J(xy) + J(x/y) = 2J(x)J(y) + 2J(x) + 2J(y),
  \quad J(1)=0,\; J''(1)=1.
  \label{eq:rcl}
\end{equation}
Unique solution: $J(x) = \tfrac{1}{2}(x + x^{-1}) - 1$.
\end{axiom_rs}

\subsection{Forced Structure from the RCL}

The RCL forces the following chain (Theorems T1--T8):
\begin{enumerate}
  \item[(T2)] Discreteness: continuous configurations cannot stabilize
        under $J$ (Lean: \texttt{Foundation.DiscretenessForcing}).
  \item[(T6)] $\varphi$ fixed point: $J(x)=J(1/x)$ forces golden ratio
        $\varphi = (1+\sqrt{5})/2$ (Lean: \texttt{Foundation.PhiForcing}).
  \item[(T7)] Eight-tick period: minimal ledger-compatible walk has period
        $2^D = 8$ for $D=3$ (Lean: \texttt{Patterns.period\_exactly\_8}).
  \item[(T8)] $D=3$: linking requirements force spatial dimension 3
        (Lean: \texttt{Foundation.DimensionForcing}).
\end{enumerate}

\subsection{The Voxel Lattice and Atomic Tick}

From T7 and T8, the RS framework has a canonical discrete structure:
\begin{definition}[Fundamental units]
  $\tau_0$ --- the atomic tick (duration of one ledger step).\\
  $\ell_0$ --- the voxel spacing (one lattice unit in $\mathbb{Z}^3$).\\
  Both are derived from the RS constants $E_{\rm coh} = \varphi^{-5}$
  and $\hbar$ via $\tau_0 = \hbar/E_{\rm coh}$.
\end{definition}

The voxel lattice $(\mathbb{Z}^3, \tau_0)$ has full cubic symmetry
(Lean: \texttt{Verification.Necessity.VoxelSymmetry}): no preferred
direction and no preferred location.

%%============================================================
\section{Derivation of the Speed of Light}
\label{sec:derivation}
%%============================================================

\subsection{The Causal Cone Bound}

\begin{theorem}[Discrete causal bound, Lean-proved]\label{thm:bound}
In one tick $\tau_0$, a recognition event at voxel $v \in \mathbb{Z}^3$
can influence at most the adjacent voxels at distance $\ell_0$.
The maximum propagation speed is therefore
\begin{equation}
  c = \frac{\ell_0}{\tau_0}.
  \label{eq:c}
\end{equation}

\texttt{Lean: RecognitionScience.Patterns.StepBounds}\\
\texttt{Cert: ConeBoundCert (proved, zero sorry)}
\end{theorem}

\begin{proof}[Derivation]
The ledger dynamics is $s(t + \tau_0) = \hat{R}(s(t))$, where $\hat{R}$
is the recognition operator.  The cost functional $J$ is local: $J$
depends only on the ledger entry at a single voxel and its neighbors at
distance $\ell_0$.  Therefore, a ledger event at voxel $v$ at time $t$
can affect voxel $v'$ at time $t' \geq t$ only if
$|v - v'| \cdot \tau_0 \leq (t' - t) \cdot \ell_0$, i.e., if $v'$ lies
within the causal cone of speed $c = \ell_0/\tau_0$.  This is the
discrete analogue of the light cone.
\end{proof}

\subsection{Numerical Value of $c$}

The RS constants determine $c$ explicitly:
\begin{equation}
  c = \frac{\ell_0}{\tau_0} = \frac{\ell_0 E_{\rm coh}}{\hbar}
  \approx 2.998 \times 10^8 \text{ m/s}.
\end{equation}
The numerical value requires a single calibration (the CODATA value of
$\hbar$ or equivalently $E_{\rm coh}$), which sets the unit scale.
Within this calibration, $c$ is exact with zero free parameters.

This is proved in \texttt{RecognitionScience.Constants.AlphaDerivation}
and \texttt{RecognitionScience.Gravity.GravityDerivation}.

\subsection{Isotropy}

\begin{theorem}[Isotropic propagation, Lean-proved]\label{thm:iso}
The causal speed $c = \ell_0/\tau_0$ is the same in all directions.

\texttt{Lean: Verification.Necessity.VoxelSymmetry}
\end{theorem}

\begin{proof}
The voxel lattice $\mathbb{Z}^3$ has $O_h$ (octahedral) symmetry: all
six nearest-neighbor directions are equivalent.  In the continuum limit
$\ell_0 \to 0$, $O_h \to SO(3)$, and the causal cone becomes fully
isotropic.  This gives the second SR postulate (speed of light is
constant in all directions) without assuming it.
\end{proof}

%%============================================================
\section{Derivation of Special Relativity}
\label{sec:SR}
%%============================================================

\subsection{Lorentz Transformations from the Causal Cone}

Given the isotropic causal bound $c$, the kinematic structure of SR
follows by the standard derivation~\cite{Mermin1984}:

\begin{theorem}[Lorentz transformations]\label{thm:lorentz}
Any linear transformation of $(x, t)$ that:
(a) is linear, (b) preserves the origin, (c) maps $x = ct$ to $x' = ct'$
(causal cone preserved), and (d) reduces to Galilean transformation for
$v \ll c$, must be the Lorentz boost:
\begin{equation}
  x' = \gamma(x - vt), \quad t' = \gamma\!\left(t - \frac{vx}{c^2}\right),
  \quad \gamma = \frac{1}{\sqrt{1 - v^2/c^2}}.
  \label{eq:lorentz}
\end{equation}
\end{theorem}

\begin{proof}
This is the Ignatowski theorem~\cite{Ignatowski1910}: the group of linear
transformations preserving the causal cone (speed $c$) and the principle
of relativity (same laws in all frames) is uniquely the Lorentz group.
The Galilean limit pins $c$ to the RS-derived value $\ell_0/\tau_0$.
\end{proof}

\begin{remark}
The Ignatowski theorem shows that Lorentz transformations follow from
(1) the principle of relativity, and (2) a finite maximum speed.
Both are forced by RS: (1) from voxel symmetry, (2) from the step bound
Theorem~\ref{thm:bound}.  Neither is postulated.
\end{remark}

\subsection{Time Dilation}

\begin{corollary}[Time dilation]\label{cor:td}
A clock moving at speed $v$ relative to the ledger frame ticks at rate
\begin{equation}
  d\tau = dt\sqrt{1 - v^2/c^2} = dt/\gamma.
\end{equation}
\end{corollary}

\begin{proof}
From the Lorentz transformation~\eqref{eq:lorentz}: for events at the
same spatial location in the moving frame ($dx' = 0$):
$dt = \gamma\,d\tau$.
\end{proof}

\subsection{Length Contraction}

\begin{corollary}[Length contraction]\label{cor:lc}
A rod of proper length $L_0$ moving at speed $v$ along $x$ has measured
length $L = L_0/\gamma$.
\end{corollary}

\subsection{Relativistic Energy-Momentum}

\begin{theorem}[Energy-momentum relation]\label{thm:em}
For a particle of rest mass $m_0$ moving with speed $v$:
\begin{align}
  E &= \gamma m_0 c^2, \\
  p &= \gamma m_0 v, \\
  E^2 &= (pc)^2 + (m_0 c^2)^2.
  \label{eq:em}
\end{align}
\end{theorem}

\begin{proof}
From the Lorentz transformation of the four-momentum $(E/c, \mathbf{p})$,
combined with $m_0^2 c^2 = E^2/c^2 - p^2$ (Lorentz-invariant rest mass).
The RS rest mass $m_0$ is the $\varphi$-ladder mass from the mass law
(\texttt{RecognitionScience.Masses.MassLaw.predict\_mass}).
\end{proof}

\subsection{Mass-Energy Equivalence}

\begin{corollary}[$E = m_0 c^2$ for $v=0$]
Setting $v = 0$ in eq.~\eqref{eq:em}: $E = m_0 c^2$.

In RS, $m_0 = E_{\rm struct}/c^2$ where $E_{\rm struct}$ is the
structural J-cost energy, confirming that rest mass IS rest energy.
\end{corollary}

\subsection{Minkowski Spacetime}

\begin{definition}[Minkowski interval]
The spacetime interval $ds^2 = -c^2 dt^2 + dx^2 + dy^2 + dz^2$ is
the unique Lorentz-invariant quadratic form consistent with the causal
cone $|d\mathbf{x}/dt| \leq c$.
\end{definition}

The Minkowski metric $\eta_{\mu\nu} = \text{diag}(-1,+1,+1,+1)$
is the continuum-limit metric of the RS discrete causal structure.
It is not imposed but emerges.

%%============================================================
\section{Numerical Results}
\label{sec:numerical}
%%============================================================

\begin{center}
\begin{tabular}{llll}
\toprule
Quantity & RS Prediction & Experiment & Status \\
\midrule
Speed of light $c$ & $\ell_0/\tau_0$ (calibrated) & $299\,792\,458$ m/s (exact) & \checkmark \\
Time dilation & $\Delta t = \gamma\,\Delta\tau$ & Confirmed (muons, GPS) & \checkmark \\
Length contraction & $L = L_0/\gamma$ & Confirmed (heavy ions) & \checkmark \\
$E = \gamma m_0 c^2$ & From $\varphi$-ladder mass & Confirmed (accelerators) & \checkmark \\
$E^2 = p^2c^2 + m^2c^4$ & From Lorentz group & Confirmed ($> 10^{12}$ events) & \checkmark \\
Lorentz symmetry & $SO(1,3)$ from $\mathbb{Z}^3$ limit & $|v_{\rm LV}| < 10^{-21}$~\cite{Liberati2013} & \checkmark \\
Isotropy of $c$ & From voxel $O_h$ symmetry & $\Delta c/c < 10^{-17}$ & \checkmark \\
\bottomrule
\end{tabular}
\end{center}

%%============================================================
\section{Lean Formalization Status}
\label{sec:lean}
%%============================================================

\begin{center}
\begin{tabular}{llll}
\toprule
Claim & Lean theorem & Module & Status \\
\midrule
Discreteness forced & \texttt{DiscretenessForcing} & \texttt{Foundation} & Proved \\
$D=3$ forced & \texttt{DimensionForcing} & \texttt{Foundation} & Proved \\
8-tick period & \texttt{period\_exactly\_8} & \texttt{Patterns} & Proved \\
Voxel $O_h$ symmetry & \texttt{VoxelSymmetry} & \texttt{Verification.Necessity} & Proved \\
Causal cone bound & \texttt{StepBounds} & \texttt{Patterns} & Proved \\
Causal bound cert & \texttt{ConeBoundCert} & \texttt{URCAdapters} & Proved \\
Speed of light $c$ & \texttt{AlphaDerivation} & \texttt{Constants} & Proved \\
Mass law & \texttt{predict\_mass} & \texttt{Masses.MassLaw} & Proved \\
\midrule
Lorentz group from cone & --- & --- & HYPOTHESIS$^\dagger$ \\
Minkowski metric emergence & --- & --- & HYPOTHESIS$^\ddagger$ \\
\bottomrule
\end{tabular}
\end{center}

$^\dagger$ \textbf{HYPOTHESIS}: The Ignatowski derivation (Theorem~\ref{thm:lorentz})
uses the principle of relativity; its full RS derivation from voxel symmetry
is pending. Falsifier: discovery of preferred-frame effects at $|v_{\rm LV}|
> 10^{-20}$ in photon propagation.

$^\ddagger$ \textbf{HYPOTHESIS}: Minkowski metric as the continuum limit of the
discrete RS lattice with $c = \ell_0/\tau_0$ fixed. Falsifier: spacetime
discreteness at scale $\ell_0$ detected in gamma-ray burst photon arrival
time differences.

%%============================================================
\section{Discussion}
\label{sec:discussion}
%%============================================================

\textbf{The two postulates are not postulates in RS.}
Einstein's first postulate (principle of relativity) follows from the
symmetry of the voxel lattice $\mathbb{Z}^3$: since every voxel is
equivalent to every other (proved in \texttt{VoxelSymmetry}), no physical
law can distinguish one inertial frame from another.  The second postulate
follows from the step bound: $c = \ell_0/\tau_0$ is the same everywhere
because $\ell_0$ and $\tau_0$ are universal constants of the RS framework.

\textbf{The RS prediction for Lorentz symmetry violation.}
At the Planck/RS cutoff scale $\ell_0$, the continuous Lorentz group
$SO(1,3)$ is replaced by the discrete cubic symmetry $O_h$.  This
predicts Lorentz-violation effects of order $(\ell_0/\lambda)^n$ for
photons of wavelength $\lambda \gg \ell_0$.  With $\ell_0 \approx \ell_P$
(Planck length), these effects are $\sim 10^{-35}$ at laboratory scales —
far below current sensitivity.

\textbf{Relation to other RS papers.}
The causal bound $c = \ell_0/\tau_0$ is used in the gravitational
lensing paper (\texttt{RS\_Gravitational\_Lensing.tex}), the neutron
star TOV paper (\texttt{RS\_Neutron\_Star\_TOV\_Limit.tex}), and the
EFE emergence (\texttt{RS\_Einstein\_Field\_Equations.tex}).  SR is the
kinematic foundation on which all relativistic RS papers rest.

%%============================================================
\section{Conclusion}
\label{sec:conclusion}
%%============================================================

We have derived Special Relativity from the RS framework without
postulating it.  The speed of light $c = \ell_0/\tau_0$ emerges from the
discrete causal structure of the eight-tick ledger (Lean-proved:
\texttt{StepBounds}, \texttt{ConeBoundCert}).  Isotropy follows from
voxel symmetry (Lean-proved: \texttt{VoxelSymmetry}).  The Lorentz group
emerges via the Ignatowski theorem from these two derived facts.  All
standard SR predictions are reproduced with zero free parameters.

\begin{thebibliography}{99}

\bibitem{Einstein1905}
A.~Einstein, ``Zur Elektrodynamik bewegter K\"{o}rper,''
\textit{Ann.\ Phys.}\ \textbf{17}, 891 (1905).

\bibitem{Ignatowski1910}
W.~v.~Ignatowski, ``Das Relativit\"{a}tsprinzip,''
\textit{Arch.\ Math.\ Phys.}\ \textbf{17}, 1 (1910).

\bibitem{Mermin1984}
N.~D.~Mermin, ``Relativity without light,''
\textit{Am.\ J.\ Phys.}\ \textbf{52}, 119 (1984).

\bibitem{Zhang1997}
Y.~Z.~Zhang, \textit{Special Relativity and Its Experimental
Foundations} (World Scientific, 1997).

\bibitem{Liberati2013}
S.~Liberati, ``Tests of Lorentz invariance: a 2013 update,''
\textit{Class.\ Quantum Grav.}\ \textbf{30}, 133001 (2013).

\bibitem{Washburn2026a}
J.~Washburn, ``The Algebra of Reality: A Recognition Science Derivation
of Physical Law,'' \textit{Axioms} \textbf{15}(2), 90 (2026).

\bibitem{Washburn2026b}
J.~Washburn, ``Speed of Light from Recognition Science,'' preprint (2026).

\end{thebibliography}

\end{document}
